\documentclass[10pt,a4paper]{article}
\usepackage[T1]{fontenc}
\usepackage[utf8]{inputenc}
\usepackage{lmodern,microtype}
\usepackage[margin=.8in]{geometry}
\usepackage{amsmath,amssymb,amsthm,booktabs,tabularx,array,enumitem,mdframed}
\usepackage{xcolor}
\usepackage[hidelinks]{hyperref}
\hypersetup{pdftitle={Regular Rotation in Sync HotStuff -- V3},pdfsubject={Agreed recovery with the original fast-path timer}}
\setlength{\parindent}{0pt}
\setlength{\parskip}{4pt}
\setlength{\emergencystretch}{2em}
\setlist{nosep,leftmargin=*}
\newcolumntype{Y}{>{\raggedright\arraybackslash}X}
\newcommand{\QC}{\mathsf{QC}}
\newcommand{\BC}{\mathsf{BC}}
\newcommand{\IC}{\mathsf{IC}}
\newcommand{\chg}[1]{\textcolor{red!70!black}{#1}}
\newtheorem{lemma}{Lemma}
\newtheorem{theorem}{Theorem}
\begin{document}
\begin{center}
{\Large\bfseries Regular Rotation in Sync HotStuff: V3}\\[3pt]
{\large Original fast-path timer; agreed recovery on failure}\\[3pt]
10 October 2026
\end{center}

\textbf{Scope and cost.} V1's first-arrival selector and V2's local
$\Delta$-window minimum fail common leader selection. V3 computes the
minimum over an agreed evidence set. It keeps immediate terminal-QC
rotation and the original vote-start $2\Delta$ timer. Its substantive
additions are a fresh recovery epoch and a failure-only authenticated
agreement on the recovery anchor and next view. The agreed anchor is
committed during recovery. This is a proof-bearing reference construction,
not a claim of globally minimal recovery communication.

\textbf{Model.} Fixed membership, $n=2f+1$, at most $f$ Byzantine replicas,
unforgeable signatures, equal-rate clocks, and known correct-to-correct
complete-message delay $\Delta$. Validation and processing are included in
this bound. Executions have finitely many events/messages in bounded time;
this is the idealized synchronous model, not a Byzantine-flooding model.
At exact boundaries, process received messages before timers. Empty batches
are available. The liveness claim is unbounded committed-height progress;
transaction inclusion additionally needs a fair request-selection policy.

\textbf{One view identity, one leader.} A view is $V=(e,r)$, ordered
lexicographically. Epoch $e$ changes only after failure recovery; normal
rotation increments $r$. The leader is the deterministic function
$L(e,r)=L_{r\bmod n}$ of this view identity, not a separately elected variable.
Every signed message binds its type, epoch, term, height and block/context
hashes. Voting identities are never reused across epochs.
Epoch $e$ starts with a common committed base $A_e$, height $H_e$, and term
$r_e$. Term $r\ge r_e$ owns fresh positions
\[
 H_e+(r-r_e)K+1,\ldots,H_e+(r-r_e+1)K,\qquad K\ge1.
\]
A normal entry requires the complete chain of intervening terminal QCs.
A QC contains $f+1$ matching votes. A quit proof is $f+1$ BLAME signatures
for one admissible $(e,r)$, or verifiable leader equivocation. Initial
entries share $A_0,r_0$ and differ in time by at most $\Delta$.

\begin{mdframed}[linewidth=.6pt,innertopmargin=5pt,innerbottommargin=5pt]
\textbf{Protocol 1. Steady state (original Figure 1 structure).}
\begin{enumerate}[label=\textbf{\arabic*.},itemsep=3pt]
\item \textbf{Propose.} The active leader proposes a valid fresh child of
its certified tip. \chg{The first child of a normal term may use the preceding
terminal QC; the first term of an epoch may extend its agreed base $A_e$.}
\item \textbf{Vote.} Process complete admissible histories even for earlier
normal terms. \chg{While unfrozen, immediately catch up through terminal QCs.
If a received proposal is ahead in height in the resulting active term,
process its missing certified ancestors in order and vote it; if already
past it on the same history, relay its QC.} Relay each proposal voted on,
multicast the vote, and start its $2\Delta$ commit timer once. Do not double-vote
at a slot. Assemble and relay QCs independently of the leader.
\chg{A terminal QC advances to $(e,r+1)$ immediately, disarms that completed
term's missing-progress listener, and preserves its pending commit timers.}
\item \textbf{Commit (non-blocking).} An uncancelled timer commits its block
and ancestors. \chg{Any admissible same-epoch quit/fork proof invokes Protocol 2
and cancels all epoch timers.} No QC-at-expiry guard is added. Never revoke
a commit. Never vote for a history conflicting with a committed prefix.
\end{enumerate}
\end{mdframed}

\begin{mdframed}[linewidth=.6pt,innertopmargin=5pt,innerbottommargin=5pt]
\textbf{Protocol 2. View change (original Figure 2 structure).}
\begin{enumerate}[label=\textbf{\roman*.},itemsep=3pt]
\item \textbf{Blame and quit.} At $s_i(e,r)+(2p+4)\Delta$, send BLAME
if fewer than $p$ valid fresh proposals are known, for $1\le p\le K$,
and no complaint is pending. \chg{On any admissible current-epoch quit proof,
relay it, freeze the whole epoch at $a_i$, and cancel all its commit timers.
Continue receiving/relaying old data and assembling QCs.}
\item \textbf{Status.} Wait $\Delta$. \chg{Snapshot all known current-epoch
QCs and quit proofs, including complete admission ancestry. Use the
failure-only agreement in Section~\ref{sec:agreement} to obtain a common
vector containing every correct snapshot.}
\item \textbf{New-view.} \chg{From that common vector select the smallest
admissible blamed term $b$ and the highest admissible current-epoch QC
(or $A_e$ if none). Let its block be $A_{e+1}$. Commit this common anchor,
and decide $V'=(e+1,b+1)$, whose leader is $L(V')$.}
\item \textbf{First vote.} \chg{Install $A_{e+1}$ and $r_{e+1}=b+1$, unfreeze,
set the new local entry clock, and resume Protocol 1. The new leader may
propose immediately. The agreed anchor needs no re-certification vote.
Retain prior records, but old-epoch proofs cannot restart this recovery.}
\end{enumerate}
\end{mdframed}

\clearpage
\section*{Validation and the recovery agreement}\label{sec:agreement}
\textbf{Complete histories and evidence.} An admissible proposal/QC has valid
application data, the scheduled leader's signature, a certified parent
(except the agreed epoch base), the prescribed slot, and a complete chain
from $A_e$. All entry QCs and signed proposals needed to verify this chain
travel with it or are already available. A same-epoch ancestry fork has a
first differing position. Canonical term/height allocation assigns that
position to one leader, giving two conflicting signatures at the same slot.
Correct nodes retain and compare all admissible histories, including those
from earlier terms, and immediately relay such evidence and freeze.

Voting compatibility and evidence validity are different checks. A quit
proof is checked against the \emph{common} epoch base and admission rules,
not against a replica's latest local commit or preferred branch. Incompatible
valid ancestry is evidence, not a reason to discard a quit proof. Old epoch
messages are archival only; future epoch messages wait for the corresponding
local recovery decision. No new data votes are made while frozen. Old votes
are still processed and can complete certificates before the snapshot.

\textbf{Local clocks.} Set $s_i(e,r)$ on entry. An entry clears the one
pending timeout complaint. Certified catch-up counts the revealed valid
proposals and disables missing-progress complaints for completed terms.
Equivocation listening remains active throughout the epoch. Recovery starts
only once per epoch; later proofs are collected without restarting timers.

\textbf{A concrete agreement implementation.} Let $a=\min_i a_i$ over correct
nodes. Immediate proof relay gives $a_i\in[a,a+\Delta]$. At
$u_i=a_i+\Delta$, replica $i$ snapshots its STATUS. A valid STATUS contains
at least one admissible current-epoch trigger and its sender's reported QC
set; the shared validation rule ignores malformed reports. Run one parallel
authenticated broadcast per sender, with domain $(\text{recovery},e,\text{sender})$.
Use the signature-chain rule of Dolev--Strong~\cite{ds}, with the following
explicit timing adaptation for start skew:
\begin{enumerate}[label=\arabic*.,itemsep=3pt]
\item Each sender signs its STATUS and sends it at its own $u_i$. Round $k$
ends locally at $u_i+2k\Delta$, for $k=1,\ldots,f+1$.
\item A chain of length $k$ contains $k$ distinct signatures, starts with
that instance's designated sender, and binds each extension to the preceding
chain. Accept it only in local round $k$; buffer early chains and discard
late ones. Record the value, subject to the common STATUS validation rule.
\item For $k\le f$, append one's signature to every accepted chain not
already containing it, and send the extended chain at the start of round
$k+1$. At the last deadline output the unique accepted valid value, or
$\bot$ if zero or multiple distinct values were accepted.
\end{enumerate}
No global clock reading or trusted timestamp is required. These are
failure-only communication rounds; normal production remains event driven.

\begin{lemma}[Common recovery vector and bounded entry skew]
Every correct node outputs the same vector, including the exact STATUS of
every correct sender. Output times differ by at most $\Delta$.
\end{lemma}
\begin{proof}
The $u_i$ differ by at most $\Delta$. A correct round-$k$ transmission is
sent at $u_i+2(k-1)\Delta$ and arrives at node $j$ by
\[
 u_i+2(k-1)\Delta+\Delta\le u_j+2k\Delta.
\]
Early arrivals are buffered, so a correct transmission is accepted in its
intended round. If a correct node accepts a value before the last round,
it either forwards an extended chain in the next round or the chain already
contains a correct signature whose prefix was broadcast. Thus all correct
nodes record that value by the final deadline. A last-round chain has
$f+1$ distinct signers, hence a correct signer; the same prefix argument
applies. Accepted value sets therefore agree. An honest designated sender
signs only its one STATUS, and everyone receives it in round 1. The common
output rule preserves that STATUS. Every node decides at $u_i+2(f+1)\Delta$,
which preserves the initial skew.
\end{proof}

\textbf{Common selection.} Collect the valid quit proofs and QCs in the
non-$\bot$ reports, including those in their admission ancestry. At least one
valid proof is present because every correct STATUS contains its trigger.
Set $b$ to the smallest term named by those proofs. The QC candidate set is
\emph{only current-epoch admissible QCs}, plus the fixed base $A_e$ as fallback;
previously hidden old-epoch QCs cannot replace that base. Break QC-rank ties
by a fixed block-hash order. Every node thus obtains the same $(b,A_{e+1})$.
Late evidence excluded from this agreed instance cannot change its output;
the next epoch rejects old-epoch triggers. This common epoch boundary
replaces V2's inconsistent local committed-view cutoff.

\clearpage
\section*{Safety proof}
\begin{lemma}[A completed fast timer protects its block]\label{lem:timer}
If a correct node votes for $B$ in epoch $e$ at time $t$ and commits it at
$t+2\Delta$, then (i) every correct node obtains $\QC(B)$ by $t+2\Delta$,
and (ii) no admissible epoch-$e$ QC conflicting with $B$ can exist, even if
assembled later from old votes.
\end{lemma}
\begin{proof}
No correct node freezes this epoch by $t+\Delta$: its relayed quit proof
would reach the committer by $t+2\Delta$ and cancel the timer. Every correct
node receives $B$ and its full admission history by $t+\Delta$. If it was
finishing the preceding recovery, the bounded epoch-entry skew ensures it
has entered epoch $e$ by then. Having already entered a later epoch would
require an earlier freeze, which has just been excluded.

A recipient on an incompatible admissible history detects the first
same-slot fork, relays the proof and freezes by $t+\Delta$, also impossible.
Otherwise it catches up through the supplied terminal certificates. If it
has already passed $B$ in height or term, its comparable certified ancestry
contains $\QC(B)$; it relays this QC. If any correct node does so, every
correct node receives the QC by $t+2\Delta$. If none has passed $B$, all
correct nodes have voted, or now vote, for $B$ by $t+\Delta$. Their at least
$f+1$ votes reach everyone by $t+2\Delta$. These votes remain processable
even if a freeze occurs after they were sent. This proves (i).

Suppose a conflicting admissible QC exists. At least one of its votes is
honest. After that voter processes $B$, seeing the conflicting history would
make it freeze instead of vote. Hence the conflicting vote must precede
its processing of $B$, no later than $t+\Delta$. Its complete proposal was
relayed at voting, and reaches the committer by $t+2\Delta$. Together with
$B$ this exposes a signed fork and cancels the timer, a contradiction.
This argument also covers a certificate assembled only after the commit.
\end{proof}

\begin{lemma}[Recovery inherits every fast commit]\label{lem:inherit}
The anchor selected at the end of an epoch extends every correct fast
commit from that epoch and the common committed base $A_e$.
\end{lemma}
\begin{proof}
Let $a$ be the first correct freeze time. For any completed timer starting
at $t$, Lemma~\ref{lem:timer}'s argument gives $a>t+\Delta$. All correct
snapshots occur at $a_i+\Delta\ge a+\Delta>t+2\Delta$, so each includes
$\QC(B)$. Agreement includes every correct snapshot. The selected QC is
therefore at least as highly ranked as $\QC(B)$ and cannot conflict with
$B$. In one epoch, canonical fresh positions increase with term and height,
so a comparable QC of higher rank extends $B$, rather than being an ancestor.
All candidates extend $A_e$ by validation. This proves the claim for every
fast commit, including ancestors committed indirectly.
\end{proof}

\begin{theorem}[Safety and common recovery leader]
All correct committed logs are prefix-comparable. For each recovery,
all correct nodes decide the same next view $(e+1,b+1)$ and hence the same
next leader, although their decision times may differ by at most $\Delta$.
\end{theorem}
\begin{proof}
Proceed by induction over epochs, beginning with the common genesis base.
Within one epoch, Lemma~\ref{lem:timer} gives a QC for every direct fast
commit and excludes an incompatible QC for any other fast commit. Thus its
fast commits are mutually compatible. Lemma~\ref{lem:inherit} proves that
the common recovery anchor extends them all. Every old timer has been
cancelled before recovery agreement finishes, so no old-epoch commit can
appear after that anchor is installed. Agreement commits one common anchor;
new-epoch histories must extend it. Fresh epoch-bound signing identities
exclude the old-view re-entry problem. This proves the induction.
The common vector, deterministic minimum and tie-break rule also give a
single next view and its deterministic leader. Temporary differences in
which term is being processed are not different recovery decisions.
\end{proof}

\textbf{Why recovery commits the anchor.} Global cancellation alone can
starve finality: with $K=1$ and fast messages, honest leaders can finish their
terms before $2\Delta$, after which a faulty successor equivocates and cancels
all pending timers. Repeating this can keep height certificates advancing
without any timed commit. A recovery decision therefore commits its common
safe anchor. This is an added rule, justified by Lemma~\ref{lem:inherit}, not
a consequence of merely re-certifying an anchor in the original protocol.

\clearpage
\section*{Liveness, iteration record, and limits}
\begin{theorem}[Committed-height progress]
Under the stated synchrony and processing assumptions, correct replicas
keep committing blocks of increasing height.
\end{theorem}
\begin{proof}
A triggered recovery terminates in finite time by the authenticated
agreement. It cannot be restarted by duplicate or late same-epoch proofs.
If an epoch continues forever without freezing, all correct replicas
receive the same admissible proposal history or detect a fork and freeze.
If progress stops, a correct missing-progress complaint is eventually
joined by the other correct replicas, whose $f+1$ signatures form a quit
proof. Otherwise the continuing proposals obtain QCs and their uncancelled
$2\Delta$ timers give unbounded commits. Thus an infinite nonfreezing epoch
cannot prevent progress.

Now suppose infinitely many recoveries occur. Every correct STATUS retains
all relevant QCs, including QCs used in normal admission and embedded in
proof histories. If a new fresh-block QC is known to an honest node before
its snapshot, the selected highest QC has height above $H_e$, and recovery
commits a taller anchor. A QC exposed only in a Byzantine report also yields
progress when it is selected. Thus infinitely many such recoveries imply
unbounded committed height.

For contradiction suppose height eventually stops increasing. Consider a
subsequent recovery without height progress. No honest node can have
normally entered a term above $r_e$: that admission requires a terminal QC
for fresh blocks, included in its STATUS and hence in the common vector.
No valid quit proof for a higher term can appear either: a BLAME certificate
contains an honest signer that entered it, and equivocation evidence must
carry its complete terminal-QC admission. Therefore every eligible quit
proof targets $r_e$, so $b=r_e$ and the next start is $r_e+1$.

Consecutive no-progress recoveries thus visit consecutive scheduled leaders.
After at most $f$ faulty starting leaders an honest one is reached. Starting
from the common anchor, its proposals reach each correct node in time to
satisfy the original $(2p+4)\Delta$ progress bounds: the leader's entry is
at most $\Delta$ later, and each proposal/QC iteration takes at most
$2\Delta$. It does not equivocate; Byzantine shares alone cannot blame it.
It produces a fresh QC. If no recovery interrupts its timer, a fast commit
occurs. If a faulty later leader triggers recovery, the fresh QC is retained
and the common anchor advances. Either outcome contradicts permanently
stopped height.
\end{proof}

\textbf{Versioned design--audit--patch record.}
\begin{tabularx}{\linewidth}{@{}p{.1\linewidth}p{.37\linewidth}Y@{}}
\toprule
Version & Selection rule & Finding or repair\\\midrule
V1 & First complete quit proof received locally. & Different nodes can
select different successors; re-entry and pending-proof handling are incomplete.\\
V2 & Wait $\Delta$, then take minimum blamed view above a local committed
cutoff. & Fails even with identical cutoffs: withheld signatures can complete
an older certificate inside only the later node's window.\\
V3 & Minimum over a common authenticated recovery vector; fresh epoch. &
The two protocols above, with safety and height-progress proofs. Normal
rotation and vote-start commit timing remain responsive.\\\bottomrule
\end{tabularx}

\textbf{V2's concrete witness.} With $n=5,f=2$, Caro's valid Alice BLAME
share exists in Round 0, but the two Byzantine shares are withheld. Alice's
terminal $X$ is nevertheless certified, and all correct nodes commit it
by Round 3; every cutoff is Alice's view 0. Bob stays silent. $P$ obtains
$\BC_1$ in Round 6, closes its window in Round 7 and chooses Caro. Caro
receives $\BC_1$ only in Round 7. The old $\BC_0$ is completed at Caro in
Round 7, after $P$'s deadline, and reaches $P$ in Round 8 after Caro's
window closes. Caro chooses Bob. Both $0$ and $1$ pass the common cutoff 0;
all correct relays take at most one round. This refutes selection agreement,
not by itself a conflicting-commit execution.

\textbf{Costs and limits.} Normal handoff has no extra waiting phase:
terminal QC immediately enables the successor. Normal data/vote forwarding
and vote-start $2\Delta$ timers are retained. Recovery now uses $n$ parallel
authenticated broadcasts of $f+1$ rounds, each padded to $2\Delta$ to absorb
entry skew. All correct nodes finish within
$(2f+4)\Delta$ of the first freeze, including the snapshot wait and start
skew. This reference agreement is deliberately not communication-optimized.
The proof does not establish that these recovery costs are necessary for
this particular selection problem, or that V3 is globally minimal. It does
not automatically prove a stake-weighted or mobile-sluggish variant.

\begin{thebibliography}{2}\small
\bibitem{sync} I. Abraham, D. Malkhi, K. Nayak, L. Ren, and M. Yin.
\emph{Sync HotStuff: Simple and Practical Synchronous State Machine Replication}.
IEEE S\&P, 2020, Section II, Figures 1--2.
\href{https://conferences.computer.org/sp/pdfs/sp/2020/349700a106.pdf}{Published paper}.
\bibitem{ds} D. Dolev and H. R. Strong.
\emph{Authenticated Algorithms for Byzantine Agreement}.
SIAM Journal on Computing 12(4), 656--666, 1983.
\href{https://doi.org/10.1137/0212045}{DOI}.
\end{thebibliography}
\end{document}
